\documentclass[a4paper,12pt]{report}

\usepackage[T2A]{fontenc}
\usepackage[utf8]{inputenc}
\usepackage[english,russian]{babel}
\usepackage{lmodern}
\DeclareFontFamily{T2A}{lmr}{}
\DeclareFontShape{T2A}{lmr}{m}{n}{<->ssub*cmr/m/n}{}
\DeclareFontShape{T2A}{lmr}{bx}{n}{<->ssub*cmr/bx/n}{}
\DeclareFontShape{T2A}{lmr}{b}{n}{<->ssub*cmr/bx/n}{}
\DeclareFontShape{T2A}{lmr}{m}{it}{<->ssub*cmr/m/it}{}
\DeclareFontShape{T2A}{lmr}{m}{sl}{<->ssub*cmr/m/sl}{}
\usepackage{graphicx}
\usepackage{setspace}
\usepackage{fancyvrb}
\usepackage{appendix}
\usepackage{listings}
\usepackage{color}
\usepackage{amssymb}
\usepackage{amsmath}
\usepackage{hyphenat}
\usepackage{subcaption}
\usepackage{hyperref}
\usepackage{float}
\usepackage{booktabs}
\usepackage{array}
\makeatletter
\renewcommand\@biblabel[1]{#1.}
\makeatother

\newcommand{\myrule}[1]{\rule{#1}{0.4pt}}
\newcommand{\sign}[2][~]{{\small\myrule{#2}\\[-0.7em]\makebox[#2]{\it #1}}}

% Поля
\usepackage[top=20mm, left=30mm, right=15mm, bottom=20mm, includefoot]{geometry}
\usepackage{indentfirst}
\footskip 20pt

\definecolor{gray}{rgb}{0.5,0.5,0.5}
\setcounter{secnumdepth}{4}

\bibliographystyle{gost780s}  %% стилевой файл для оформления по ГОСТу

\lstset{
	basicstyle=\footnotesize,           % the size of the fonts that are used for the code
	numbers=left,                   % where to put the line-numbers
	numberstyle=\tiny\color{gray},
	captionpos=b,                   % sets the caption-position to bottom
	breaklines=true,                % sets automatic line breaking
}


\begin{document}
	
	%%%%%%%%%%%%%%%%%%%%%%%%%%%%%%%
	%%%                         %%%
	%%% Начало титульного листа %%%
	
	\thispagestyle{empty}
	{ \large
		%%% Название работы %%%
		\begin{center}
			\renewcommand{\baselinestretch}{1.2}
			\Large
			{\Large Анализ зашумленных и частично наблюдаемых данных судовых навигационных систем}
		\end{center}
		\medskip%
		
		\vspace{3\bigskipamount}%
		\begin{flushright}
			\begin{minipage}{0.5\textwidth}
				\begin{flushleft}
					Гашков С.П.
				\end{flushleft}
			\end{minipage}
		\end{flushright}
		
		\vfill%
	}
	
	%%% Конец титульного листа  %%%
	%%%                         %%%
	%%%%%%%%%%%%%%%%%%%%%%%%%%%%%%%
	
	% Размер шрифта (14пт)
	\large
	
	% Межстрочный интервал (полуторный)
	\onehalfspacing
	
\clearpage
\renewcommand{\thesection}{\arabic{section}}
\graphicspath{{../pict/}}

\section{Постановка}

По точкам \((t_i,\mathrm{lat}_i,\mathrm{lon}_i)\) восстанавливается бинарная разметка судового трека. Метка \(y_i=0\) — аномалия, \(y_i=1\) — штатный ход. Точки стоянки в выборку не входят. Положительный класс всех метрик ниже — аномалия.

Файл содержит \(419\,453\) строк. Фрагмент индексов \([200\,000,\,350\,000)\) --- \(150\,000\) строк. Удалены \(6\,842\) точек стоянки; пустых координат нет. Остаётся \(143\,158\) точек, упорядоченных по времени. На этом фрагменте сделаны два разреза.

Разрез 1: обучение \(43\,158\) точек (аномалий \(62{,}0\,\%\)), тест \(100\,000\) (аномалий \(53{,}9\,\%\), \(19\) событий), \(L_{\mathrm{ref}}=2181{,}7\) км. Разрез 2: обучение \(93\,158\) (аномалий \(45{,}1\,\%\)), тест \(50\,000\) (аномалий \(77{,}2\,\%\), \(7\) событий), \(L_{\mathrm{ref}}=560{,}4\) км. Тест разреза 2 --- последние \(50\,000\) точек теста разреза 1. Карты ниже показывают это окно; метрики разреза 1 посчитаны по всем \(100\,000\) точкам.

Модель даёт \(s_i\): большее значение — более аномальная точка. Порог выбирается на всём обучении. Сначала максимизируется
\[
\mathrm{MCC}=\frac{TP\cdot TN-FP\cdot FN}{\sqrt{(TP+FP)(TP+FN)(TN+FP)(TN+FN)}},
\]
затем среди порогов в полосе \(0{,}02\) от лучшего MCC берётся минимум ошибки длины. Разметка \(\hat a_i=\mathbf{1}[s_i>\tau]\), \(\hat y_i=0\) при аномалии. Исключения без такого порога — DBSCAN, One-Class SVM и MPCT.

\section{Метрики}

Расстояние — гаверсинус сферы \(R=6371\) км:
\[
d_{\mathrm{geo}}(p,q)=2R\arcsin\sqrt{\sin^2\frac{\Delta\varphi}{2}+\cos\varphi_1\cos\varphi_2\sin^2\frac{\Delta\lambda}{2}}.
\]
Хорда разрыва линейна по широте и по кратчайшей дуге долготы. Время в \(d_{\mathrm{geo}}\) не входит.

Точечные полнота, точность и \(F_1\) считают каждую точку отдельно, без времени и координат:
\[
\mathrm{Precision}=\frac{TP}{TP+FP},\quad
\mathrm{Recall}=\frac{TP}{TP+FN},\quad
F_1=\frac{2TP}{2TP+FP+FN}.
\]
При нулевом знаменателе доля не определена. \(F_1=0\) при \(TP=0\) и \(FP+FN>0\); при \(TP=FP=FN=0\) не определена и \(F_1\).

PR-AUC — площадь под кривой precision--recall по score. Порог кривой задаёт \(\hat a_i(\tau)=\mathbf{1}[s_i\ge\tau]\); узлы — различные \(s_i\) по убыванию, кривая дополняется точкой \((0,1)\), интеграл берётся трапециями. Без аномалий в \(G\) или без score величина не определена.

Affiliation \(F_1\) сравнивает аномальные интервалы по времени. Участок \(\{i:y_i=0\}\) становится полуинтервалом \([t_{\mathrm{start}},t_{\mathrm{end}+1})\). Ось режется на зоны ближайшего истинного события. Локальные \(p_j\) и \(r_j\) — survival function временного расстояния внутри зоны. Precision усредняется по зонам \(S\), куда попало предсказание, recall — по всем \(m\) событиям:
\[
\mathrm{Precision}=\frac{1}{|S|}\sum_{j\in S}p_j,\quad
\mathrm{Recall}=\frac{1}{m}\sum_{j=1}^{m}r_j,\quad
F_1=\frac{2\,\mathrm{Precision}\cdot\mathrm{Recall}}{\mathrm{Precision}+\mathrm{Recall}}.
\]
При \(S=\emptyset\) precision и \(F_1\) не определены. Координаты не используются.

Задержка события — время от его начала до первой точки с \(\hat y=0\):
\[
\mathrm{Delay}_k=t_k^{\mathrm{detect}}-t_k^{\mathrm{start}}.
\]
Пропуск в среднее не входит и нулём не заменяется. Где ниже названа задержка, это среднее по обнаруженным событиям. На разрезе 1 медиана равна нулю у всех моделей, нашедших хотя бы одно событие. На разрезе 2 она нулевая не у всех: у Robust PCA \(15\) с, у LSTM и DONUT \(625\) с, у GP \(60\) с, у TCN \(5\) с.

Пространственные метрики берут координаты \(G\). В сумму входят только пропуски \(S_{FN}=\{i:y_i=0\land\hat y_i\neq 0\}\), для которых у разрыва есть оба штатных соседа (множество \(V\)):
\[
RMSE_d=\sqrt{\frac{1}{|S_{FN}\cap V|}\sum_{i\in S_{FN}\cap V}d_i^2},\qquad
H_{FN}=\max_{i\in S_{FN}\cap V}d_i.
\]
При пустом \(S_{FN}\) обе величины равны \(0\). Симметричный Хаусдорф двух множеств не считается.

Distance loss сравнивает длины ломаных:
\[
DL=|L_{\mathrm{pred}}-L_{\mathrm{ref}}|,\qquad DLR=\frac{DL}{L_{\mathrm{ref}}}.
\]
\(L_{\mathrm{ref}}\) — штатные точки \(G\) и хорды внутренних разрывов. \(L_{\mathrm{pred}}\) — точки, оставленные штатными. При \(L_{\mathrm{ref}}=0\) отношение \(DLR\) не определено.

\section{Модели}

Кроме трёх фильтров состояния, точка описывается шагом \((\Delta x_i,\Delta y_i)\), скоростью, ускорением, поворотом курса и формой окна \(w=6\):
\[
\mathrm{shape}_i=\frac{\sum_{k=i-w}^{i-1}d_{\mathrm{geo}}(p_k,p_{k+1})}{\max\bigl(d_{\mathrm{geo}}(p_{i-w},p_i),1\bigr)}\quad(i>w).
\]
Калман, фильтр частиц и GP-EKF ведут \((E,N,v_E,v_N)\) в равнопромежуточной проекции от медианы обучающего хода. DBSCAN, LOF и RPCA считают плотность или разложение того ряда, который размечают. Isolation Forest и HMM видят обучающую выборку с обеими метками. Остальные учатся на штатном ходе. Оконные сети обучены \(6\) эпох.

\begin{table}[H]
\centering
\caption{Правило детекции. \(s_i\) растёт с аномальностью; \(\hat a_i=\mathbf{1}[s_i>\tau]\), если не сказано иное.}
\footnotesize
\singlespacing
\begin{tabular}{@{}>{\raggedright\arraybackslash}p{0.22\textwidth}>{\raggedright\arraybackslash}p{0.74\textwidth}@{}}
\toprule
Модель & Правило \\
\midrule
ARIMA\((2,1,1)\) & Две модели широты и долготы на гладком участке хода. \(s_i=d_{\mathrm{geo}}(z_i,\hat z_i)\). \\
Калман & Постоянная скорость, \(s_i=r_i^{\top}S_i^{-1}r_i\). Шаг \(\le 45\) с; после разрыва \(>180\) с состояние сбрасывается. \\
Particle filter & Та же кинематика, \(K=32\). \(s_i=-\log\bigl(\frac1K\sum_j p(z_i\mid x_i^{(j)})\bigr)\). \\
GP & Два RBF-процесса шага, \(250\) точек хода, \(\ell=6\) ч. Сумма квадратов нормированных остатков шага. \\
GP-EKF & Постоянная скорость и GP-остаток перехода по скорости и \(\Delta t\), \(160\) опор. \(s_i\) --- NIS. \\
HMM & \(o_i=(v_i,a_i^{\mathrm{acc}},\Delta\psi_i)\). Разность предсказательных логарифмов аномальной и штатной модели. \\
Hampel & Скорость, окно \(21\) точка. \(s_i=|v_i-m_i|/(1{,}4826\,\mathrm{MAD}_i)\). \\
DBSCAN & Шаг и \(\log\mathrm{shape}\), \(\mathrm{MinPts}=8\). На разрезе 1 \(\varepsilon=0{,}86\). На разрезе 2 подбор не завершён. \\
LOF & \(s_i=\mathrm{LOF}_{35}\) по шагу, скорости, ускорению и повороту. \\
Isolation Forest & Лес по всей обучающей выборке. \(s_i=2^{-E[h]/c(n)}-\frac12\). \\
One-Class SVM & RBF, \(\nu=0{,}1\), \(4\,000\) точек хода. Аномалия при \(s_i>0\). \\
Robust PCA & Окна длины \(16\). \(s_i=\max_{W\ni i}\|S_W(i)\|\). \\
LSTM, DONUT & Окно относительных координат, \(s_i=\max_{W\ni i}s_W\). LSTM --- диагональный Махаланобис, DONUT --- ошибка восстановления и \(D_{\mathrm{KL}}\). \\
TCN & Каузальный прогноз следующей точки. Диагональный Махаланобис остатка. \\
Anomaly Transformer & \(\mathrm{Softmax}(-\mathrm{AssDis})\odot\|x-\hat x\|_2^2\), обучение только ошибкой восстановления. \\
STGVAD & Окно одной траектории, ребро при расстоянии \(<5\) км и только из прошлого. Ошибка прогноза положения. \\
MPCT & Якорь --- последняя принятая точка. Разрез 1: \(v_{\max}=8\) м/с. Разрез 2: \(v_{\max}=10\) м/с. Форма \(\le 1{,}5\). \\
\bottomrule
\end{tabular}
\end{table}

\section{Результаты}

Индекс 1 --- тест из \(100\,000\) точек, индекс 2 --- из \(50\,000\). Столбцы \(n_1\) и \(n_2\) --- число найденных событий из \(19\) и из \(7\). Прочерк --- подбор DBSCAN на разрезе 2 оборвался нехваткой памяти.

\begin{table}[H]
\centering
\caption{Точечные метрики двух разрезов.}
\footnotesize
\singlespacing
\setlength{\tabcolsep}{4pt}
\begin{tabular}{@{}lrrrrrr@{}}
\toprule
Модель & MCC\(_1\) & \(F_{1,1}\) & \(n_1\) & MCC\(_2\) & \(F_{1,2}\) & \(n_2\) \\
\midrule
DBSCAN & 0,823 & 0,92 & 19 & --- & --- & --- \\
Robust PCA & 0,738 & 0,88 & 18 & 0,690 & 0,90 & 6 \\
MPCT & 0,736 & 0,87 & 18 & 0,644 & 0,89 & 7 \\
Isolation Forest & 0,662 & 0,85 & 19 & 0,856 & 0,97 & 7 \\
Калман & 0,646 & 0,80 & 18 & 0,572 & 0,84 & 7 \\
HMM & 0,560 & 0,76 & 18 & 0,654 & 0,88 & 7 \\
One-Class SVM & 0,527 & 0,80 & 19 & 0,706 & 0,92 & 7 \\
ARIMA\((2,1,1)\) & 0,526 & 0,67 & 18 & 0,426 & 0,69 & 7 \\
Anomaly Transformer & 0,495 & 0,79 & 18 & 0,768 & 0,94 & 6 \\
GP-EKF & 0,376 & 0,68 & 16 & 0,356 & 0,71 & 7 \\
Particle filter & 0,365 & 0,50 & 19 & 0,293 & 0,52 & 7 \\
STGVAD & 0,285 & 0,70 & 18 & 0,506 & 0,81 & 6 \\
DONUT & 0,243 & 0,64 & 16 & 0,483 & 0,75 & 6 \\
LSTM & 0,235 & 0,61 & 16 & 0,431 & 0,72 & 6 \\
Hampel & 0,146 & 0,26 & 17 & 0,119 & 0,25 & 7 \\
LOF & 0,073 & 0,19 & 17 & 0,003 & 0,21 & 7 \\
TCN & 0,062 & 0,49 & 17 & 0,384 & 0,63 & 6 \\
GP & $-0{,}109$ & 0,69 & 19 & 0,081 & 0,06 & 7 \\
Все штатные & 0 & 0 & 0 & 0 & 0 & 0 \\
Все аномалии & 0 & 0,70 & 19 & 0 & 0,87 & 7 \\
\bottomrule
\end{tabular}
\end{table}

\begin{table}[H]
\centering
\caption{Длина и геометрия пропуска. \(H_{FN}\) --- в километрах. «Ход» --- доля сохранённых штатных точек, \%. \(L_{\mathrm{ref}}\) равен \(2181{,}7\) км и \(560{,}4\) км.}
\footnotesize
\singlespacing
\setlength{\tabcolsep}{3.5pt}
\begin{tabular}{@{}lrrrrrr@{}}
\toprule
Модель & DLR\(_1\) & \(H_1\) & Ход\(_1\) & DLR\(_2\) & \(H_2\) & Ход\(_2\) \\
\midrule
DBSCAN & 9,11 & 926 & 82,7 & --- & --- & --- \\
Robust PCA & 80,24 & 13\,058 & 88,6 & 160,27 & 7\,497 & 95,0 \\
MPCT & 1,35 & 382 & 90,2 & 3,39 & 23 & 91,4 \\
Isolation Forest & 5,12 & 81 & 64,5 & 27,27 & 49 & 92,0 \\
Калман & 7,42 & 459 & 93,2 & 16,73 & 459 & 92,9 \\
HMM & 32,37 & 1\,000 & 88,2 & 81,11 & 460 & 94,6 \\
One-Class SVM & 17,65 & 926 & 61,4 & 56,74 & 459 & 90,7 \\
ARIMA\((2,1,1)\) & 26,41 & 1\,000 & 95,8 & 74,19 & 460 & 96,6 \\
Anomaly Transformer & 76,61 & 13\,058 & 56,8 & 157,46 & 7\,497 & 92,4 \\
GP-EKF & 31,59 & 1\,039 & 74,5 & 69,38 & 461 & 85,4 \\
Particle filter & 40,88 & 1\,031 & 95,0 & 121,17 & 574 & 95,7 \\
STGVAD & 78,36 & 13\,058 & 52,7 & 147,09 & 7\,497 & 90,6 \\
DONUT & 51,48 & 13\,058 & 62,7 & 90,33 & 7\,497 & 96,7 \\
LSTM & 52,38 & 13\,058 & 67,7 & 93,69 & 7\,497 & 93,4 \\
Hampel & 49,22 & 2\,664 & 93,3 & 138,02 & 461 & 94,9 \\
LOF & 51,09 & 1\,288 & 93,3 & 133,19 & 574 & 88,4 \\
TCN & 79,62 & 13\,058 & 62,9 & 168,15 & 7\,497 & 98,0 \\
GP & 22,77 & 1\,000 & 0,0 & 150,52 & 574 & 100 \\
Все штатные & 96,34 & 13\,058 & 100 & 215,67 & 7\,497 & 100 \\
Все аномалии & 1,00 & 0 & 0 & 1,00 & 0 & 0 \\
\bottomrule
\end{tabular}
\end{table}

\(L_{\mathrm{ref}}\) второго разреза меньше в \(3{,}89\) раза, поэтому тот же абсолютный \(DL\) дал бы \(DLR\) в \(3{,}89\) раза больше. Маска «все аномалии» на втором тесте имеет \(F_1=0{,}87\): это доля аномалий, а не качество детектора.

На разрезе 1 точечный лидер --- DBSCAN: \(\mathrm{MCC}=0{,}82\), полнота \(0{,}98\), все \(19\) событий, средняя задержка \(6\) с. Кластеры строятся по тестовому облаку, на обучении MCC равен \(0{,}42\). \(DLR=9{,}1\): на карте ход вдоль Волги режется ложными тревогами. Маршрут держит MPCT, \(DLR=1{,}35\), ход \(90\,\%\). Robust PCA рядом по \(F_1\) и оставляет пропуск \(H_{FN}=13\,058\) км, общий с маской без вырезов, LSTM, TCN, DONUT, STGVAD и Anomaly Transformer. Калман сохраняет \(93\,\%\) хода при полноте \(0{,}71\) и \(DLR=7{,}4\). GP-EKF на обучении имеет MCC \(0{,}78\), на тесте \(0{,}38\): ниже Калмана при той же кинематике.

На разрезе 2 лидер по точкам --- Isolation Forest: \(\mathrm{MCC}=0{,}86\), \(F_1=0{,}97\), ход \(92\,\%\) вместо \(65\,\%\) на первом тесте, \(H_{FN}=49\) км, \(DLR=27\). MPCT при \(v_{\max}=10\) м/с снова держит ломаную: \(DLR=3{,}39\), \(H_{FN}=23\) км, ход \(91\,\%\). Падение MCC у MPCT, Калмана, ARIMA и фильтра частиц близко к сдвигу доли аномалий с \(54\,\%\) до \(77\,\%\). Рост MCC у леса, SVM, HMM и оконных сетей --- в основном снятые ложные тревоги. Полнота LSTM, TCN, DONUT и STGVAD не выросла, общий дальний пропуск стал \(7\,497\) км. GP-EKF снова ниже Калмана: \(0{,}36\) против \(0{,}57\). У GP медиана score хода выше медианы аномалий на обоих разрезах; порог перешёл от \(0{,}002\) к \(8{,}1\), и модель сохраняет весь ход при полноте \(0{,}03\).

\begin{figure}[H]
\centering
\includegraphics[width=\textwidth]{v1_truth.png}
\caption{Окно карт: последние \(50\,000\) точек разреза 1, они же весь тест разреза 2. Синий --- штатный ход, красный --- аномалия, оранжевый пунктир --- хорда эталона. Подложка --- OpenStreetMap.}
\end{figure}

\begin{figure}[H]
\centering
\includegraphics[width=\textwidth]{v1_dbscan.png}
\caption{DBSCAN, разрез 1, то же окно. Синий --- верно оставленный ход, зелёный --- найденная аномалия, красный крест --- пропуск, оранжевый --- ложная тревога, чёрная линия --- принятый маршрут. В легенде рисунка буква TP стоит на штатной точке, а TN --- на найденной аномалии. В метриках отчёта положительный класс --- аномалия, поэтому зелёная точка есть TP.}
\end{figure}

\begin{figure}[H]
\centering
\includegraphics[width=\textwidth]{v2_iforest.png}
\caption{Isolation Forest, разрез 2, то же окно. Обозначения те же, что на рис.~2.}
\end{figure}

\begin{figure}[H]
\centering
\includegraphics[width=\textwidth]{v1_mpct.png}
\caption{MPCT, разрез 1, \(v_{\max}=8\) м/с и форма \(\le 1{,}5\). Обозначения те же, что на рис.~2.}
\end{figure}

\begin{figure}[H]
\centering
\includegraphics[width=\textwidth]{v2_mpct.png}
\caption{MPCT, разрез 2, \(v_{\max}=10\) м/с и форма \(\le 1{,}5\). Обозначения те же, что на рис.~2.}
\end{figure}

\section{Выводы}

На обоих разрезах \(F_1\) и \(DLR\) упорядочивают модели по-разному. Высокий \(F_1\) совместим с непригодной ломаной. У Robust PCA и Anomaly Transformer он около \(0{,}9\), а \(DLR\) составляет десятки и сотни эталонов. MPCT --- единственная модель, которая оба раза держит ошибку длины на порядке одного эталона и около \(90\,\%\) хода. Для восстановления маршрута он работоспособен. Порог скорости на сетке сдвинулся с \(8\) до \(10\) м/с, форма окна осталась \(1{,}5\): правило устойчиво, рабочую точку нужно выбирать заново. Это основная перспектива среди проверенных детекторов маршрута.

Точечные лидеры разрезов не совпадают. Первый --- DBSCAN, и его MCC частично описывает тестовое облако. На удлинённом обучении расчёт DBSCAN не завершён. На втором впереди лес. Он работоспособен как метка точки, когда обучение длиннее и менее аномально, и на первом разрезе даёт кратчайший геометрический хвост, \(81\) км. По длине он остаётся позади MPCT.

Калман на обоих разрезах сохраняет ход и точность выше \(0{,}9\) при полноте около \(0{,}7\). Фильтр частиц при той же кинематике держит полноту около \(0{,}35\). GP-EKF, в котором гауссовский процесс учит остаток перехода, на обучении выглядит сильно и на обоих тестах проигрывает Калману. Перспективна постоянная скорость как консервативный фильтр. Остаток GP и частицы на этом треке развивать не стоит: они не добавляют ни полноты, ни короткой ломаной. ARIMA устойчиво точна и сохраняет ход, но забирает лишь около половины аномальных точек.

Оконные сети, Robust PCA и Anomaly Transformer на обоих разрезах оставляют один дальний скачок принятого маршрута. Дополнительный штатный ход снимает у них ложные тревоги и не поднимает полноту. Пока обучение штрафует только ошибку окна, эта группа не пригодна для маршрута. Hampel, LOF и GP класс не отделяют ни на одном разрезе. Подбор порога на этих признаках не перспективен.

\end{document}
